\section{Beurling-Helson und Folgerungen}

\begin{genericthm}{Beurling-Helson}\label{skript:1.1}
Sei $E$ ein $\mathscr{S}$-invarianter Unterraum von $\L^2$.
Dann gibt es folgende zwei Fälle:
\begin{enumerate}[label=\textbf{(\arabic*)}]
\item
Es gilt $\mathscr{S} E = E$.
Dann existiert eine $\lambda_N$-messbares $e \subset \mathbb{T}$ mit $E = \chi_e \L^2$.
\item
Es gilt $\mathscr{S} E \neq E$.
Dann existiert ein $\lambda_N$-messbares $\Theta : \mathbb{T} \to \C$ mit $|\Theta| = 1$ fast überall, sodass
$E = \Theta \H^2$ gilt.
\end{enumerate}
\end{genericthm}

\begin{proof}
\begin{enumerate}[label=\textbf{(\arabic*)}]
\item
Sei $P_E$ die orthogonale Projektion auf $E$ und $\varphi := P_E 1 $.
Das ist $1- \varphi$ wegen $\varphi E = E$ orthogonal zu $\mathscr{S}^n \varphi$ für alle $n \in \Z$.
Aufgrund von 
\begin{align*}
\langle \mathscr{S}^n \varphi, 1 - \varphi \rangle 
= \int \limits_\mathbb{T} \mathscr{S}^n \varphi \overline{(1 - \varphi)} \td{\lambda_N}
= \int \limits_\mathbb{T} (\varphi - | \varphi|^2) z^n \td{\lambda_N}
= 0
\end{align*}
erhalten wir $\varphi = | \varphi|^2$ und finden ein $e \subset \mathbb{T}$ mit $\varphi = \chi_e$.
Da durch $\lbrace z^n \ | \ n \in \Z \rbrace $ ein Basis von $\L^2$ gegeben ist, erhalten wir mit
\begin{align*}
\bigvee_{n \in \Z} \mathscr{S}^n \chi_e :=
\overline{\lin} \bigcup_{n \in \Z} \lbrace \mathscr{S}^n \chi_e \rbrace
= \chi_e \cdot \L^2,
\end{align*}
dass $\chi_e \L^2 \subset E$ gilt.
Wir müssen also noch die Gleichheit zeigen.
Hierfür wählen wir $\alpha \in E \ominus \chi_e \L^2 $.
Damit gilt 
\begin{align*}
\langle 1- \chi_e , \mathscr{S}^n \alpha \rangle = 
\int \limits_\mathbb{T} (1- \chi_e) \overline{\alpha} z^{-n} \td{\lambda_N} = 0
\end{align*}
für $n \in \Z$. Damit folgt $(1- \chi_e) \overline{\alpha} = 0$.
Analog folgt, dass 
\begin{align*}
\forall_{n \in \Z}  \ : \ \langle \alpha , \mathscr{S}^n \chi_e \rangle = 0
\quad \Rightarrow \quad \overline{\alpha} \chi _e= 0
\end{align*}
gilt. Wegen 
\begin{align*}
(1- \chi_e) \overline{\alpha} =
\overline{\alpha} \chi _e= 0
\end{align*}
erhalten wir direkt $\alpha = 0$.
Damit gilt $E = \chi_e \L^2$.

\item
In diesem Fall ist $\mathscr{S}E$ ein echter Teilraum von $E$.
Wir wählen ein $\Theta \in E \ominus \mathscr{S} E$ mit $\| \Theta \|_2 = 1$.
Der Operator $\mathscr{S}$ ist unitär, weswegen wir
\begin{align*}
\langle \mathscr{S}^n \Theta , \Theta \rangle
= 
\langle \Theta, \mathscr{S}^{-n} \Theta \rangle
=
\int \limits_\mathbb{T} z^n | \Theta |^2 \td{\lambda_N} =0 
\end{align*}
für $n \in \N$ erhalten.
Damit ist $|\Theta| $ konstant.
Da $\lambda_N(\mathbb{T}) = 1$ gilt, folgt mit 
\begin{align*}
\langle \Theta, \Theta \rangle 
=
\int \limits_\mathbb{T} | \Theta|^2 \td{\lambda_N} = 1 
\end{align*}
auch $| \Theta  | = 1 $ fast überall.
Wir betrachten $\Theta$ nun als Multiplikationsoperator, d.h. $f \mapsto \Theta f$. 
Damit ist $\Theta$ eine Isometrie auf $\L^2$.
Mit der Einbettung
\begin{align*}
\iota \ : \H^2 \hookrightarrow \L^2, \quad 
z^n \mapsto z^n
\end{align*}
und der Isometrie erhalten wir durch
\begin{align*}
\bigvee_{n \in \N_0} \mathscr{S}^n \Theta
=
\Theta \bigvee_{n \in \N_0} \mathscr{S}^n 1 
=
\Theta \bigvee_{n \in \N_0} \iota(z^n) 1
= \Theta H^2, 
\end{align*}
dass $\Theta \H^2 \subset E$ gilt.
Übrig zu zeigen ist wieder die Gleichheit.
Wir wählen $\alpha \in E \ominus \Theta \H^2$.
Dann gilt 
\begin{align*}
\langle \alpha , \mathscr{S}^n \Theta \rangle
= 
\langle \mathscr{S}^{-n} \alpha , \Theta \rangle
= 0
\end{align*} 
für $n \in \N_0$ und wegen $ \mathscr{S}^n \alpha \in \mathscr{S} E$ gilt auch
\begin{align*}
\langle \Theta, \mathscr{S}^n \alpha \rangle 
=
\langle \mathscr{S}^{-n} \Theta, \alpha \rangle
= 0
\end{align*}
für $n \in \N$.
Damit erhalten wir
\begin{align*}
\int \limits_\mathbb{T} \Theta \overline{\alpha} z^n \td{\lambda_N} = 0
\end{align*}
für alle $n \in \Z$.
Damit folgt direkt $\alpha = 0$.
Somit gilt $E = \Theta H^2$.
\end{enumerate}
\end{proof} 

\begin{lemma}\label{skript:1.2}
Sei $E$ ein Unterraum von $\L^2$.
Es gelten:
\begin{enumerate}[label=\textbf{(\arabic*)}]
\item
Jeder Unterraum der Form $\chi_e \L^2$ ist $\mathscr{S}$ und $\mathscr{S}^{-1}$-invariant.
Die Darstellung $E = \chi_e \L^2$ ist eindeutig.

\item
Jeder Unterraum der Form $\Theta H^2$, $| \Theta | = 1$ f.ü. ist $\mathscr{S}$-invariant, jedoch nicht $\mathscr{S}^{-1}$-invariant.
Die Darstellung $E = \Theta H^2$ ist bis auf eine multiplikative Konstante eindeutig. 

\end{enumerate}
\end{lemma}

\begin{proof}
\begin{enumerate}[label=\textbf{(\arabic*)}]
\item
Folgt direkt aus dem letzten Beweis.
\item
Die Invarianz bzgl. $\mathscr{S}$ und Nichtinvarianz bzgl. $\mathscr{S}^{-1}$ folgen auch hier direkt aus dem letzten Beweis.
Wir betrachten nun $\Theta \H^2 = \Theta^\prime \H^2$.
Dann folgt $\Theta \overline{\Theta^\prime} \H^2 = H^2$
bzw $\H^2 = \Theta^\prime \overline{\Theta} \H^2$.
Nun gilt $\Theta \overline{\Theta^\prime} \in \H^2$ und $\Theta \overline{\Theta^\prime} \in \overline{\H^2}$.
Seien $f , \overline{f} \in \H^2$.
Dann erhalten wir 
\begin{align*}
\hat{f}(n) = 0
\end{align*}
und
\begin{align*}
\overline{\hat{f}}(n) = \overline{\hat{f}(-n)} = 0
\end{align*}
für $n < 0$. Somit ist $f$ konstant und wir erhalten die Eindeutigkeit bis auf eine multiplikative Konstante.
%Mit
%\begin{align*}
%\bigvee_{n\in \N_0} \mathscr{S}^n \Theta 
%= 
%\bigvee_{n \in \N_0} \mathscr{S}^n \Theta^\prime
%\quad \Rightarrow \quad
%\Theta \mathscr{S}^n 1 = C \cdot \Theta^\prime \mathscr{S}^n 1
%\end{align*}
%folgt die restliche Aussage sofort. 
\end{enumerate}
\end{proof}

\begin{genericthm}{Folgerung}\label{skript:1.3}
Sei $E \subset \H^2$ ein nicht trivialer Unterraum und $\mathrm{S}E \subset E$.
Dann gilt $E = \Theta \H^2$ mit $| \Theta | = 1$ fast überall auf $\mathbb{T}$.
\end{genericthm}

\begin{proof}
Mit der Einbettung 
\begin{align*}
\iota \ : \H^2 \hookrightarrow \L^2, \quad 
z^n \mapsto z^n
\end{align*}
gilt $\iota(E) \subset \iota(\H^2) \subset \iota(\L^2)$ und $\mathscr{S} E  \subset E$.
Aus Beurling-Helson erhalten wir sofort
\begin{align*}
E = \iota(E) = \Theta \H^2
\end{align*}
mit $|\theta | = 1$ f.ü. auf $\mathbb{T}$.
Wegen $E \subset \H^2$ gilt auch $\Theta \in \H^2$.
\end{proof}


\begin{genericthm}{Folgerung}\label{skript:1.4}
Sei $f \in \H^2$ und $\lambda_N( \lbrace \xi \in \mathbb{T}\ | \ f(\xi) = 0 \rbrace) > 0$.
Dann gilt $f = 0$.
\end{genericthm}

\begin{proof}
Zunächst ist klar, dass
\begin{align*}
A:=\lbrace \xi \in \mathbb{T}\ | \ f(\xi) = 0 \rbrace
\end{align*}
bis auf eine Nullmenge eindeutig ist. Wir betrachten nun $f$ als Element des $\L^2$.
Dann ist 
\begin{align*}
E_f = \bigvee_{n \in \N_0} \mathscr{S}^n f
\end{align*}
eine $\mathscr{S}$-invarianter Unterraum.
Die in $E_f$ enthaltenen Funktionen verschwinden auf $A$.
Aus $f \neq 0$ folgt zunächst $E_f \neq {0}$.
Mit \ref{skript:1.3} erhalten wir 
$E_f = \Theta \H^2$ mit $\Theta \in E_f$ und $|\Theta | = 1$ fast überall.
\end{proof}